\documentclass[10pt, a4paper]{article}

% --- Packages ---
\usepackage[utf8]{inputenc}
\usepackage[margin=1.5cm]{geometry} % Extremely tight margins
\usepackage{graphicx}
\usepackage{caption}
\usepackage{enumitem}
\usepackage{mathpazo} % Professional, compact font

\begin{document}

% --- Compact Header ---
\begin{center}
    \Large\textbf{Expected Key Figure Submission}\\[0.2cm]
    \normalsize\textbf{Topic 7:} Cost-Aware Machine Learning for Accelerated Drug Discovery\\
    \textbf{Students:} [Your Name] \& [Partner's Name] \quad \textbf{Supervisor:} Michael Backenköhler\\
    \textbf{Date:} May 15, 2026
\end{center}
\vspace{-0.4cm} % Pull the image up

% --- The Figure ---
\begin{figure}[h!]
    \centering
    % strictly limits image height to leave room for text
    \includegraphics[height=6.5cm, keepaspectratio]{Honest_Expected_Key_Figure.png}
    \vspace{-0.2cm}
    \caption{\textbf{Expected Cost-Accuracy Pareto optimization.} \textit{(Note: Because baseline models are undergoing hyperparameter tuning to resolve overfitting, the y-axis values are simulated projections based on an initial Ridge Regression baseline of 0.416).} The plot demonstrates the intended visualization of the trade-off between computational cost and predictive accuracy. The green star marks the expected optimal sparse pipeline identified via Cost-Weighted Lasso.}
    \label{fig:pareto}
\end{figure}

\vspace{-0.6cm} % Pull the text up closer to the image

% --- The 6 Required Questions ---
\section*{Detailed Figure Description}
\vspace{-0.2cm}

\noindent\textbf{1. What is the key finding you want to communicate?}\\
An exhaustive virtual screening workflow is unnecessary. By applying a Cost-Weighted $L_1$ penalty, we can identify a ``Pareto Optimal'' pipeline of computationally cheap docking algorithms that achieves near-maximum accuracy while drastically reducing overall computational cost compared to using exhaustive tools alone.

\vspace{0.2cm}
\noindent\textbf{2. What is on the x-axis?}\\
Total Computational Cost per molecule (in seconds). Because the cost discrepancy between tools is massive (0.31s for CNNscore vs. 407.5s for DiffDock), this axis uses a logarithmic scale to clearly visualize both fast and slow algorithms.

\vspace{0.2cm}
\noindent\textbf{3. What is on the y-axis?}\\
Predictive Accuracy, quantified by the Pearson correlation coefficient ($r$), which measures how closely the ML model's predicted binding affinity matches the actual lab-tested biological affinity (\texttt{true\_value}).

\vspace{0.2cm}
\noindent\textbf{4. What does the legend show?}
\begin{itemize}[noitemsep, topsep=0pt, leftmargin=*]
    \item \textbf{Grey Dots:} Individual docking algorithms and sub-optimal scoring pipelines.
    \item \textbf{Red Dashed Line:} The Pareto Frontier connecting the non-dominated subsets of tools.
    \item \textbf{Green Star:} Our proposed ``Optimal Pipeline'' selected by the Lasso regularizer.
\end{itemize}

\vspace{0.2cm}
\noindent\textbf{5. What is the caption?}\\
Provided directly beneath Figure \ref{fig:pareto}. It explains the trade-off, highlights the optimal pipeline, and includes a disclaimer regarding the projected nature of the current accuracy values.

\vspace{0.2cm}
\noindent\textbf{6. What data processing is needed to create it?}\\
Raw docking logs are aggregated by molecule ID (taking \texttt{min}/\texttt{max} based on the tool's physical metric). Failed poses receive worst-case imputation. Outliers (e.g., HYDE engine failures) are mitigated via IQR clipping and Robust Scaling. Finally, data is fed into a Cost-Weighted Lasso regression, which iteratively forces the weights of expensive, low-impact tools to zero.

\end{document}